% PNP REVIEW REPORT::REDTEAM::RT-0014::2026-10-11T16:36:00Z
\documentclass[11pt]{article}
\usepackage[margin=0.84in]{geometry}
\usepackage{amsmath,amssymb,amsthm,booktabs,enumitem,xurl,hyperref}
\setlength{\emergencystretch}{3em}
\newtheorem{definition}{Definition}
\newtheorem{theorem}{Theorem}
\newtheorem{corollary}[theorem]{Corollary}
\newcommand{\EXACT}{\textbf{PROVED/EXACT}}
\newcommand{\KNOWN}{\textbf{KNOWN THEOREM}}
\newcommand{\OPEN}{\textbf{OPEN}}
\newcommand{\RETIRED}{\textbf{FALSIFIED/RETIRED}}
\newcommand{\COMP}{\textbf{COMPUTATIONAL}}
\title{UNIVERSAL\_CORE\_REDTEAM RT-0014:\\
Forward Child Locality Does Not Supply Reverse Fanout}
\author{P versus NP project review}
\date{11 October 2026, 16:36 UTC}
\begin{document}
\maketitle

\begin{abstract}
TESTS-0009 correctly showed that a bounded-fanout circuit with a forward
child accessor admits a linear-size, bounded-occurrence proper-3-CNF
compilation without parent queries.  It then left source-specific
\texttt{JMV-FORWARD-FANOUT1} open.  This report isolates the exact remaining
gap.  In an acyclic child-only circuit oracle with one target gate and $n$
potential later parents, deciding whether the target has fanout zero,
enumerating its complete parent list, or certifying fanout at most a fixed
$\Delta<n$ requires $n$ child queries in the deterministic worst case.
This is the elementary OR adversary, stated in the correct circuit
realization.  Therefore the generic implication ``forward child locality
supplies a usable fanout certificate'' is false; a repair must prove
additional structure of the particular Jahanjou--Miles--Viola labels or
retain TESTS-0009's parent-owned clause construction.  The result is only an
input-oracle query lower bound.  It is not a lower bound for arbitrary SAT
deciders or global CNF compilers.  P versus NP remains open.
\end{abstract}

\section{Executive status and exact target}

\begin{center}
\small
\begin{tabular}{@{}p{0.43\linewidth}p{0.19\linewidth}p{0.30\linewidth}@{}}
\toprule
Claim & Status & Verdict\\
\midrule
\texttt{FORWARD-FANOUT-DETECTION1} & \EXACT, pending independent review &
$n$ child queries are necessary and sufficient in the declared oracle model.\\
\texttt{CHILD-LOCALITY-IMPLIES-FANOUT-CERT1} & \RETIRED &
Forward access alone contains no sublinear reverse-incidence certificate.\\
Finite audit & \COMP & 376,370 vector/order executions through $n=12$; zero failures.\\
\texttt{JMV-FORWARD-FANOUT1} & \OPEN &
The source-specific label theorem has not been unfolded.\\
Arbitrary-SAT transfer, UCORE-TIME1, P versus NP & \OPEN &
No unrestricted computation lower bound follows.\\
\bottomrule
\end{tabular}
\end{center}

The attacked generic candidate is
\[
 \forall C\ \bigl[\text{gate labels expose type and children locally}\bigr]
 \Longrightarrow
 \bigl[\text{fanout and complete parents are locally available}\bigr].
\tag{CF}
\]
The smallest source-specific repair would replace the antecedent by the
actual JMV label grammar and explicitly construct either (i) a bounded
visible fanout theorem, or (ii) a locally addressable copy tree.  No such
theorem is claimed here.

\section{Child-only circuit oracle}

\begin{definition}[Acyclic forward-incidence family]
Fix an input gate $u$, private dummy input gates $d_1,\ldots,d_n$, and
later unary gates $p_1,\ldots,p_n$.  For $z\in\{0,1\}^n$, the only child of
$p_i$ is $u$ when $z_i=1$ and $d_i$ when $z_i=0$.  A forward query at
$p_i$ returns that child label.  Thus
\[
 \operatorname{fanout}(u)=|z|_1,
 \qquad
 \operatorname{Parents}(u)=\{p_i:z_i=1\}.
\]
The graph is acyclic because all $p_i$ follow their children.  The encoded
input length is $N=\Theta(n\log n)$ under explicit indexed labels; the
number of oracle positions is exactly $n$.
\end{definition}

\begin{theorem}[\texttt{FORWARD-FANOUT-DETECTION1}]
Let $A$ be any deterministic adaptive algorithm receiving the labels and
forward child oracle of the preceding family.  Each of the following tasks
has worst-case child-query complexity exactly $n$:
\begin{enumerate}[nosep]
\item decide whether $\operatorname{fanout}(u)=0$;
\item output the complete set $\operatorname{Parents}(u)$;
\item for fixed $0\le\Delta<n$, decide whether
      $\operatorname{fanout}(u)\le\Delta$.
\end{enumerate}
Equivalently the lower bound is $\Omega(N/\log N)$ in encoded length.
\end{theorem}

\begin{proof}
For task 1, answer every query with its private dummy.  If $A$ stops after
fewer than $n$ queries, compare the all-zero completion with the completion
having one at an unqueried position.  The complete transcript is identical
but the correct answers differ.  Hence $A$ must query all $n$ positions on
the all-zero input.  The same pair shows that a claimed parent list cannot
be certified complete before all positions are read.

For task 3, answer one on the first $\Delta$ queried positions and zero on
later queries.  If $A$ stops before querying every position, the unqueried
bits can be completed either all zero, of weight $\Delta$, or with one
additional one, of weight $\Delta+1$.  Again the transcript is identical and
the predicate differs.  Scanning all positions gives the matching upper
bound for all three tasks.
\end{proof}

\paragraph{Accounting.}
An offline inverse index is not a counterexample: constructing it performs
the same $n$ child queries and stores the reverse-incidence information
(up to $O(n\log n)$ bits for explicit parent labels).  Later constant-time
lookups are charged only after that preprocessing.  Nonuniform advice can
also contain the parent set, but then the advice is the missing resource and
must be included in SAME\_REALIZATION.  A one-bit output does not reveal the
parent list.

\section{What this does and does not repair}

TESTS-0009's forward-wire compilation remains valid.  If every gate has
constant fanin, constant fanout, fixed gate templates, and a child accessor,
then each parent may own the clauses describing its outgoing wires.  No
target gate has to enumerate its parents.  The theorem above therefore does
\emph{not} refute that construction; it refutes only an attempted upgrade
from child locality to a reverse-port certificate.

Conversely, a normalizer that must allocate all leaves of a target-owned
copy tree, emit a complete target-owned incident-clause block, or certify
that the source fanout is at most $\Delta$ inherits the $n$-query bound in
this interface.  The bound can be evaded only by an additional source-label
invariant, a supplied inverse index, or a different global construction.
This is the precise missing arrow:
\[
 \text{JMV output-bit locality}
 \not\Rightarrow_{\rm generic}
 \text{bounded visible fanout or local reverse incidence}.
\]

The JMV theorem states a local clause-generation result: an explicit circuit
produces an indexed clause, with each output bit depending on only constantly
many index bits in the strongest regime \cite{jmv}.  It does not, by that
statement alone, expose the complete set of labels pointing to a chosen
gate.  A source-specific proof may still exist and remains \OPEN.

\section{SAME\_REALIZATION and hostile controls}

The oracle task takes a target label and must compute a reverse-incidence
predicate.  SAT takes an encoded formula and emits one satisfiability bit.
These are different functions.  No compiler from an arbitrary uniform
deterministic SAT decider to the child-only reverse-incidence task is known.
Accordingly, the theorem cannot be charged to every SAT algorithm.

The exact indistinguishability pair survives permutation of labels, clause
order, alternate encodings that do not add a reverse port, independent
products, padding, redundant clauses, and DIRECT/COPY/DELAY/INTERP/RECOMP
wrappers.  DELAY adds charged work but no information.  Horn, 2SAT,
affine/Tseitin Gaussian elimination, and bounded-treewidth solvers may decide
their formulas without reconstructing this parent set; they are controls
against, not support for, a universal transfer.  A parent-owned global
compiler is another valid hostile control.  Semantic fanout complexity is
therefore not identified with actual SAT cost.

\section{Bounded reproducible experiment}

The dependency-free invocation
\begin{verbatim}
timeout 30s python3 experiments/review/forward_fanout_detection_redteam.py
\end{verbatim}
completed in $0.597$ seconds; bytecode compilation passed.  For each
$1\le n\le12$, the program enumerated all $2^n$ incidence vectors and up to
50 fixed query orders.  Across 376,370 vector/order executions, exact
zero-fanout scanning always had worst case $n$, while complete enumeration
always used $n$.  Every one of 2,664 capped shortcuts had an unqueried
singleton adversary.  Thirty-two bounded-fanout transcript witnesses and
eight hostile-control classes also passed.  These observations are
\COMP; the theorem above is symbolic for every $n$.

The script SHA-256 is
\nolinkurl{ad0d10e9e1ae670551aa4bbdc61f644d5ba2cfd35649f6c0d09feb4129f1d956};
the result SHA-256 is
\nolinkurl{27fc48977d51b905a5038300a15b884b57b578944b03d849aa33a84a0803a465}.

\section{Barrier and cross-domain transfer audit}

\begin{description}[leftmargin=0cm,labelwidth=0cm]
\item[] \textbf{DIRECTLY TRANSFERABLE THEOREM.}
The OR/weight decision-tree adversary transfers exactly to the declared
child-oracle family, because one oracle query reveals one incidence bit.
Offline reverse indexing transfers exactly only when its full scan and
storage are charged in the same realization.
\item[] \textbf{USEFUL ANALOGY / CANDIDATE.}
Reverse-adjacency data structures, VLSI fanout/routing, distributed in-neighbor
discovery, observability, and LWE-style external extraction all suggest that
missing inverse access must be paid for.  A new model-matched compiler is
needed before any of them constrains an arbitrary SAT decider.
\item[] \textbf{MISLEADING / FAILED TRANSFER.}
Circuit fanout, routing congestion, communication cuts, proof-system size,
Schur-square growth, quadratic relations, Hilbert functions, syzygies,
collision entropy, special soundness, and random-oracle binding do not by
themselves lower-bound unrestricted one-bit SAT computation.
\end{description}

The query proof relativizes and arithmetizes without difficulty; this is a
warning, not evidence of barrier circumvention.  Baker--Gill--Solovay show
that relativizing methods cannot settle P versus NP by themselves
\cite{bgs}, and Aaronson--Wigderson show that algebrization is a further
barrier \cite{aw}.  Natural Proofs is not the immediate hinge because no
large constructive property against general circuits is proved
\cite{rr}.  Arbitrary polynomial exponents remain untouched: an
$\Omega(N/\log N)$ restricted query bound is polynomial and compatible with
SAT in P.  Assuming SAT in P supplies a polynomial fallback to every
recognizer-plus-fallback wrapper and yields no contradiction here.

\paragraph{Post-quantum ledger.}
The LWE worst-case/average-case architecture is an ANALOGY/CANDIDATE for
external black-box extraction; the missing solver-to-oracle compiler is
exactly the unproved bridge.  Schur-square/componentwise products, quadratic
relation spaces, higher Hilbert invariants, collision extraction and
collapse-binding are FAILED TRANSFERS for this inverse-incidence obstacle.
No cryptographic assumption is admitted.

\section{Corpus-wide reconciliation}

\subsection*{CORPUS REVIEWED}
The run surveyed the full repository tree and read root \texttt{main.tex},
actual \texttt{main4.tex}/\texttt{main5.tex}, \texttt{AGENTS.md},
\texttt{CODEX.md}, the complete pnp-research skill, red-team mode and
protocol, all current task specifications TASK-001--007, the newest verified
historical MAIN, current status and coordination surfaces, governance notes,
and the HEC/VSC-HEC/G-PPC-HEC/CRS-PPC/CCSPC causal branch.  Materially
relevant reports included TESTS-0009, MODELBRIDGE-0014, LB-0014,
SYNTHESIS-0009 and RT-0013 with their cited artifacts.  The predecessor
TESTS-0009 exists in its canonical GitHub directory, is substantive,
contains its expected header and final sentinel, and was read back.
Primary literature on local reductions and the relativization,
algebrization and Natural Proofs barriers was checked before promotion.

\subsection*{DUPLICATION / CONTRADICTION AUDIT}
MODELBRIDGE-0014 already proved that a supplied parent port is sufficient and
gave a permutation-oracle separation between child and parent access.
RT-0013 separated forward tags from inverse incidence.  TESTS-0009 then
showed that parent-owned clauses bypass parent enumeration.  The present
result does not repeat those conclusions: it gives the exact deterministic
$n$-query theorem for fanout-zero, complete enumeration, and bounded-fanout
certification in an acyclic child-oracle circuit family.  It is consistent
with, and narrows, TESTS-0009's open JMV handoff.

\subsection*{WHAT CHANGED RELATIVE TO ROOT MAIN}
Root MAIN correctly records the forward child port as insufficient for
generic reverse indexing while leaving source-specific JMV structure open.
The new exact theorem upgrades that warning to a tight query bound and
retires the unqualified fanout-certificate implication.  It does not change
the executive OPEN verdict.  This review role has no authorized root-MAIN
ownership region, so root \texttt{main.tex} is not edited; main4/main5
receive review includes for synthesis arbitration.

\subsection*{ADOPTED / REJECTED / RETIRED PRIOR COMPONENTS}
Adopted are TESTS-0009's parent-owned forward-wire construction,
MODELBRIDGE-0014's interface separation, RT-0013's inverse-incidence audit,
and LB-0014's rule that optional annotations cannot be charged to arbitrary
deciders.  Rejected is any silent inference from local child labels to a
local reverse port.  Retired is
\texttt{CHILD-LOCALITY-IMPLIES-FANOUT-CERT1} as an unqualified interface
claim.  Source-specific JMV fanout and local copy-tree theorems are not
retired.

\subsection*{STALE OR OVERSTATED CLAIMS FOUND}
Any wording that calls the JMV open arrow merely ``mechanical label
unfolding'' is overstated until the label grammar is shown to bound or
locally enumerate incoming references.  Conversely, wording that says a
parent port is necessary for every bounded-occurrence compilation is stale:
TESTS-0009's parent-owned construction is a verified counterexample.  The
correct distinction is target-owned reverse normalization versus global
parent-owned emission.

\section{Smallest repair and synthesis handoff}

The next proof obligation is source-specific and binary:
\begin{enumerate}
\item unfold the exact JMV gate-label grammar and prove that every relevant
gate has $O(1)$ locally addressable parents or a locally constructible copy
tree, including all preprocessing and label lengths; or
\item exhibit a JMV label family realizing the child-only adversary, and
retain parent-owned clauses instead of reverse normalization.
\end{enumerate}
Even a successful first branch establishes a local CNF compiler, not an
arbitrary-SAT lower bound.  The separate load-bearing arrow remains an
external, uniform, model-matched transformation from every hypothetical
polynomial SAT decider to a realization that must solve an independently
hard resource problem.  No such transformation is known.

\begin{thebibliography}{9}
\bibitem{jmv} H.~Jahanjou, E.~Miles and E.~Viola,
\emph{Local reductions}, ICALP 2015, pp.~749--760.
\url{https://arxiv.org/abs/1311.3171}.
\bibitem{bgs} T.~Baker, J.~Gill and R.~Solovay,
\emph{Relativizations of the P=?NP question}, SIAM Journal on Computing
4(4) (1975), 431--442.
\bibitem{aw} S.~Aaronson and A.~Wigderson,
\emph{Algebrization: A new barrier in complexity theory}, ACM Transactions
on Computation Theory 1(1) (2009), Article 2.
\url{https://www.scottaaronson.com/papers/alg.pdf}.
\bibitem{rr} A.~A.~Razborov and S.~Rudich,
\emph{Natural proofs}, Journal of Computer and System Sciences 55(1)
(1997), 24--35.
\end{thebibliography}

\paragraph{Final status.}
\texttt{FORWARD-FANOUT-DETECTION1}: \EXACT, pending independent review.\\
Finite audit: \COMP.\\
\texttt{CHILD-LOCALITY-IMPLIES-FANOUT-CERT1}: \RETIRED.\\
\texttt{JMV-FORWARD-FANOUT1}, arbitrary-SAT transfer, UCORE-TIME1,
nonrelativization, nonalgebrization and P versus NP: \OPEN.
P versus NP remains open.

% END_OF_VERIFIED_REPORT::REDTEAM::RT-0014
\end{document}
